\section{Institutional background}\label{sec:background}

\subsection{Program mechanics}

The NMTC, enacted in the Community Renewal Tax Relief Act of 2000, awards
investors a 39\% federal tax credit, claimed over seven years, on Qualified
Equity Investments (QEIs) made into certified Community Development
Entities. CDEs in turn deploy that capital as Qualified Low-Income
Community Investments into Qualified Active Low-Income Community
Businesses (QALICBs) located in eligible census tracts. The public release
distinguishes four QALICB types: real-estate projects (RE), non-real-estate
operating businesses (NRE), special-purpose entities (SPE), and
CDE-to-CDE loans. These mechanics put the CDE at the center of the program.
It holds the allocation and originates the deals, then structures the
capital stack around the subsidized tranche. Figure~\ref{fig:mechanics}
shows the sequence. Capital moves from the investor through the CDE to the
project, while the tax credit returns to the investor.

\begin{figure}[t]
\centering
\resizebox{\textwidth}{!}{% The program, as a working sketch. Fineliner boxes, fountain-pen money
% flows, pencil annotations. Drawn to be read in ten seconds: capital walks
% left to right, the credit walks back, and the CDE sits where every
% decision happens.
\begin{tikzpicture}[x=1cm, y=1cm, font=\small]

% ── actors ──
\node[sketchbox, minimum width=2.5cm, minimum height=1.5cm, align=center]
  (inv) at (0,0) {\textbf{Investor}\\[-1pt]{\footnotesize private capital}};
\node[washbox, minimum width=2.9cm, minimum height=1.8cm, align=center]
  (cde) at (4.6,0) {\textbf{CDE}\\[-1pt]{\footnotesize certified intermediary}\\[-1pt]{\footnotesize \emph{holds the allocation}}};
\node[sketchbox, minimum width=2.7cm, minimum height=1.5cm, align=center]
  (qal) at (9.4,0) {\textbf{QALICB}\\[-1pt]{\footnotesize the project}};

% ── the eligible-tract ground the project stands in ──
\draw[pencilline] (7.7,-1.28) .. controls (8.4,-1.5) and (10.4,-1.5) .. (11.1,-1.28);
\node[sketchnote, anchor=north] at (9.4,-1.42)
  {low-income tract: poverty $\geq 20\%$ or MFI $\leq 80\%$};

% ── flows forward ──
\draw[inkarrow] (inv.east) -- node[above, sketchlabel] {QEI} (cde.west);
\draw[inkarrow] (cde.east) -- node[above, sketchlabel] {QLICI} (qal.west);

% ── the credit walks home along the top ──
\draw[inkarrow] (cde.north) .. controls (4.6,1.75) and (0,1.75) .. (inv.north);
\node[sketchlabel, anchor=south] at (2.3,1.62)
  {39\% federal credit, over 7 years};

% ── private stack joins at the project ──
\draw[pencilarrow] (9.4,2.1) -- (qal.north);
\node[sketchnote, anchor=south, align=center] at (9.4,2.18)
  {other private debt \& equity\\[-2pt] joins the stack here};

% ── the object of study ──
\draw[pencilline] (11.35,-0.55) -- (11.75,-0.55) -- (11.75,0.55) -- (11.35,0.55);
\node[sketchnote, rotate=90, anchor=south] at (12.35,0)
  {Leverage $=$ cost/QLICI};

% ── the observed ledger, underlined like a margin scribble ──
\node[sketchnote, anchor=north west] at (-1.2,-2.05)
  {the public ledger records both dollars of the ratio, project by project, 2001--2022};
\draw[pencilline] (-1.2,-2.62) -- (7.6,-2.62);

\end{tikzpicture}
}
\caption{The program, as a working sketch. Private capital enters as a
Qualified Equity Investment. The intermediary deploys it as a QLICI into a
project in an eligible tract, and the 39\% credit returns to the investor
over seven years. The public ledger records both inputs to the leverage
ratio for every project.}
\label{fig:mechanics}
\end{figure}

\subsection{Eligibility and the 20\% non-metro minimum}

Eligible low-income community (LIC) tracts are defined by a poverty rate of
at least 20\% or median family income at or below 80\% of the area
benchmark, with a targeted-populations override. Relevant here, a later
amendment directs the Treasury to write rules ensuring that
non-metropolitan counties receive a proportional allocation of qualified
equity investments. That instruction is IRC \S45D(i)(6), added by the Tax
Relief and Health Care Act of 2006 (P.L. 109--432, \S102(b)), so it
postdates the program's creation in the Community Renewal Tax Relief Act of
2000 (P.L. 106--554, \S121), leaving the earliest sample years outside its
scope. The statute names no percentage. The familiar 20\% figure comes
from the CDFI Fund's administration of the requirement through allocation
applications and agreements. I therefore use ``administrative target''
throughout.\footnote{The effective dates matter for the analysis;
Section~\ref{sec:bunching} reports the mandate results for the full period
and for origination years from 2007 onward.} Aggregate deployment is close
to that line: non-metro tracts received 19.6\% of QLICI dollars over the
sample period.

\subsection{The intermediary layer}

The public release names 345 distinct CDEs in the transaction ledger
(343 at the project level); 310 executed five or more transactions. The
twenty largest account for 23.9\% of QLICI dollars, so deployment is
moderately concentrated. Specialization is extreme: among CDEs with at
least five transactions, cumulative non-metro shares span the full range
from 0\% to 100\%, with 39\% of such CDEs never deploying non-metro and
6.5\% deploying there at least four times out of five.
This dispersion motivates the decomposition. An aggregate rural penalty
can arise when rural deals flow through intermediaries that mobilize less
\emph{everywhere}, even if their rural and urban projects have similar
leverage. The public release identifies CDEs by name. A later extension will
classify their institutional form (bank subsidiary, nonprofit CDFI,
for-profit specialist).
